\usepackage[spanish]{babel}   % español
\usepackage[utf8]{inputenc}   % acentos sin codigo
\usepackage[right=2cm,left=2.5cm,top=2cm,bottom=2cm,headsep=0.5cm,footskip=0.8cm]{geometry}
\usepackage{amssymb, amsmath} % Paquetes matemáticos de la American Mathematical Society
\usepackage{graphicx}		  % para imagenes
\usepackage{float} 		      % para imagenes en lugar solicitado H
\usepackage{listings}	      % para codigo fuente
\usepackage{enumerate}        % para listas numeradas con letras
\usepackage{natbib}			  % para referencias bibliograficas
\usepackage{url}			  % referencias bibliograficas WEB
\usepackage{ragged2e}         % alineacion de texto
\usepackage{subcaption}	      % para imagenes multiples
\usepackage{hyperref} 		  % para generar marcadores en el PDF ponerlo antes del \begin{document}
\usepackage{booktabs}         % Líneas elegantes en tablas
\usepackage{threeparttable}   % Notas al pie en tablas
\usepackage{tabularx}         % Tablas con ancho controlado (\textwidth)
\usepackage{array}            % Modificadores de columna (>{...})

%para imagenes y tablas estilo APA7
\captionsetup[figure]{labelfont={bf}, textfont={it} ,labelformat={default}, labelsep=newline, name={Figura }, justification=raggedright,singlelinecheck=false} %-----formato del caption figure
\captionsetup[table]{labelfont={bf}, textfont={it} ,labelformat={default}, labelsep=newline, name={Tabla },justification=raggedright,singlelinecheck=false} %-----formato del caption table

\hypersetup{hidelinks,
		colorlinks=true,
		allcolors=black,
		pdfstartview=Fit,
		breaklinks=true
	}	%para quitar marco rojo

% quitar el 0.
\renewcommand\thesection{\arabic{section}}